\subsection{COLLECTIVIZING RELATIONS}

Let R be a relation and let x be a letter. If y and \(y'\) denote letters distinct from x which do not appear in R, then the relations

\[
(\exists y) (\forall x) ((x \in y) \leftrightarrow R), (\exists y ^ {\prime}) (\forall x) ((x \in y ^ {\prime}) \leftrightarrow R)
\]

are identical by CS8 (Chapter I, §4, no. 1). The relation thus defined (which does not contain x) is denoted by Coll \(_{x}\) R.

If \(\operatorname{Coll}_{\mathbf{x}} R\) is a theorem in a theory \(\mathcal{G}\) , \(R\) is said to be collectivizing in \(\mathbf{x}\) in \(\mathcal{G}\) . When this is so, we may introduce an auxiliary constant \(\mathbf{a}\) , distinct from \(\mathbf{x}\) and the constants of \(\mathcal{G}\) , and which does not appear in \(\mathbf{R}\) , with the introductory axiom \((\forall \mathbf{x}) ((\mathbf{x} \in \mathbf{a}) \Longleftrightarrow \mathbf{R})\) , or equivalently (if \(\mathbf{x}\) is not a constant of \(\mathcal{G}\) ) \((\mathbf{x} \in \mathbf{a}) \Longleftrightarrow \mathbf{R}\) .

Intuitively, to say that R is collecting in x is to say that there exists a set a such that the objects x which possess the property R are precisely the elements of a.

\minorheading{Examples}

\begin{enumerate}
\def\labelenumi{(\arabic{enumi})}
\item
  The relation \(x \in y\) is clearly collectivizing in \(x\) .
\item
  The relation \(x \notin x\) is not collectivizing in \(x\) ; in other words, (not \(\operatorname{Coll}_x(x \notin x)\) ) is a theorem. Let us argue by contradiction and suppose that \(x \notin x\) is collectivizing. Let \(a\) be an auxiliary constant, distinct from \(x\) and the constants of the theory, with the introductory axiom
\end{enumerate}

\[
(\forall x) ((x \notin x) \iff (x \in a)).
\]

Then the relation

\[
(a \notin a) \iff (a \in a)
\]

is true by C30 (Chapter I, §4, no. 3). The method of disjunction of cases (Chapter I, §3, no. 3) shows first that the relation \(\pmb{a} \notin \pmb{a}\) is true, and then that the relation \(\pmb{a} \in \pmb{a}\) is true, which is absurd.

C49. Let R be a relation and x a letter. If R is collectivizing in x, the relation \((\forall x)((x \in y) \Longleftrightarrow R)\) , where y is a letter distinct from x which does not appear in R, is functional in y.

This follows immediately from C48.

Very often in what follows we shall have at our disposal a theorem of the form \(Coll_{x}R\) . To represent the term

\[
\tau_ {y} (\forall x) ((x \in y) \iff R),
\]

which does not depend on the choice of the letter y (distinct from x and not appearing in R), we shall introduce a functional symbol \(\mathcal{E}\mathbf{x}(\mathbf{R})\) ; the corresponding term does not contain X. This term is denoted by ``the set of all x such that R''. By definition (Chapter I, §4, no. 1) the relation

\[
(\forall x) ((x \in \mathcal {E} x (R)) \iff R)
\]

is identical with \(Coll_{x}R\) ; consequently the relation R is equivalent to

\[
\boldsymbol {x} \in \mathcal {E} _ {\boldsymbol {x}} (\boldsymbol {R}).
\]

C50. Let R, S be two relations and let x be a letter. If R and S are collectivizing in x, the relation \((\forall x)(R \Longrightarrow S)\) is equivalent to

\[
\mathcal {E} _ {\boldsymbol {x}} (\boldsymbol {R}) \subset \mathcal {E} _ {\boldsymbol {x}} (\boldsymbol {S}),
\]

and the relation \((\forall x)(R \Longleftrightarrow S)\) is equivalent to \(\mathcal{E}_{x}(R) = \mathcal{E}_{x}(S)\) .

This follows immediately from the preceding remark and from Definition 1 and axiom A1.

